\documentclass{article}

\usepackage{neurips_2026}

\usepackage[utf8]{inputenc}   % allow utf-8 input
\usepackage[T1]{fontenc}      % use 8-bit T1 fonts
\usepackage{hyperref}         % hyperlinks
\usepackage{url}              % simple URL typesetting
\usepackage{booktabs}         % professional-quality tables
\usepackage{amsfonts}         % blackboard math symbols
\usepackage{amsmath}          % math environments
\usepackage{amssymb}          % math symbols
\usepackage{nicefrac}         % compact symbols for 1/2, etc.
\usepackage{microtype}        % microtypography
\usepackage{xcolor}           % colors
\usepackage{graphicx}         % figures
\usepackage{tikz}             % diagrams

\title{Learning to Rotate: Temporal and Semantic Rotary Encoding for Sequential Modeling}

\author{%
  Hailing Cheng \\
  LinkedIn Inc.\\
  Mountain View, CA 94043 \\
  \texttt{haicheng@linkedin.com} \\
  \And
  Daqi Sun \\
  LinkedIn Inc.\\
  Mountain View, CA 94043 \\
  \texttt{daqsun@linkedin.com} \\
  \And
  Xinyu Lu \\
  LinkedIn Inc.\\
  Mountain View, CA 94043 \\
  \texttt{xinlu@linkedin.com} \\
}

\begin{document}

\maketitle

\begin{abstract}
Every Transformer architecture dedicates enormous capacity to learning rich representations in \emph{semantic embedding space}---yet the rotation manifold acted upon by Rotary Positional Embeddings (RoPE) has been treated as a fixed, hand-crafted structure, populated only by discrete ordinal indices. We argue that this \emph{rotation space} is a largely overlooked second dimension of expressivity in the attention mechanism, one whose systematic exploration may open a new door for attention-based architectures. The analogy to complex numbers is instructive: just as introducing the imaginary axis---orthogonal to and independent of the real line---unlocked new algebraic structure once believed impossible, treating the rotation manifold as a learnable, signal-conditioned space opens an orthogonal degree of freedom in attention. In this framing, the token embedding encodes the \emph{semantic} (real) component of a representation---\emph{what} a token means---while the rotation encodes its \emph{dynamic} (imaginary) component---\emph{how} it relates to every other token across time, position, and context.

We introduce \textbf{SIREN-RoPE}, a concrete instantiation of this idea, which populates the rotation dimension with heterogeneous signals---continuous timestamps, cyclical temporal patterns, and categorical metadata---via a dual-branch Sinusoidal Representation Network (SIREN). As a proof of concept, we evaluate on a production-scale news feed dataset from a major social network, demonstrating that activating this hidden dimension yields consistent improvements across calibration and ranking objectives with negligible computational overhead. We invite the community to view the rotation space not as a solved positional-encoding detail, but as an untapped axis whose rich structure may prove as consequential for attention as the imaginary unit proved for algebra.
\end{abstract}


%--------------------------------------------------------------------
\section{Introduction}
\label{sec:intro}
%--------------------------------------------------------------------

The meteoric rise of Transformer architectures in sequence modeling is largely attributed to the self-attention mechanism. Because self-attention is inherently permutation-invariant, the burden of maintaining structural order and modeling the causal relations falls entirely on positional encodings and causal masking. While Rotary Position Embeddings (RoPE) \citep{su2024roformer} have improved the modeling of relative distances by framing position as a planar rotation of query/key subspaces, they remain fundamentally limited by a \textit{one-size-fits-all} design: they rely on fixed inverse-frequencies and treat sequence positions as simple discrete integers.

In many high-stakes domains---sequential recommendation, event-stream modeling, and large language modeling---this discrete representation is insufficient. An interaction that occurred seven days ago is semantically distinct from one that occurred seven minutes ago, yet traditional models treat them as computationally identical if they share the same relative index. These failures, however, point to a deeper structural gap: the rotation manifold of RoPE has been studied almost exclusively as a \emph{positional scaffolding device}---a tool for translational equivariance---while its potential as a learnable, content-sensitive space remains almost entirely unexplored.

We draw an analogy to the history of complex numbers. For centuries, mathematicians worked within the real line until introducing an imaginary axis---orthogonal to and independent of the real axis---unlocked entirely new algebraic structure, from the Fundamental Theorem of Algebra to Fourier analysis and modern signal processing. The attention mechanism presents an analogous opportunity. Token embeddings constitute the \emph{real} (semantic) dimension: they encode \emph{what} a token means. The rotation applied by RoPE constitutes an orthogonal \emph{dynamic} dimension: it encodes \emph{how} that token relates to every other token in the sequence across time and structure. Treating this rotation dimension as a learnable, signal-conditioned space---rather than a fixed positional scaffold---is the analogue of introducing the imaginary axis: a conceptual extension that opens an entirely new expressive structure for attention.

\subsection{Key challenges}

Specifically, we identify four open gaps in the
literature:
\begin{enumerate}
  \item \textbf{The rotation dimension is neglected.} RoPE defines a rotation manifold over the query/key space, but this manifold has been treated as a fixed, hand-crafted structure entirely determined by ordinal position. The expressive capacity of this space---its ability to encode temporal dynamics, semantic context, or any structured relationship---remains almost entirely untapped. We argue that unlocking this dimension is as conceptually significant as extending the real line to the complex plane.
  \item \textbf{Fixed frequency rigidity.} Standard RoPE uses hand-crafted inverse-frequency schedules that cannot adapt to the statistical distribution of the input corpus.
  \item \textbf{Temporal–ordinal conflation.} A token's ordinal rank in the sequence is not a faithful proxy for its temporal recency; real-world event logs feature irregular and bursty arrival patterns.
  \item \textbf{Absence of semantic geometry.} The rotation applied to queries and keys does not vary with token content or category, missing an important opportunity to modulate attention in semantic space.
\end{enumerate}

\subsection{Our contributions.} We introduce \textbf{SIREN-RoPE} as both a demonstration and an argument: a concrete system that shows what becomes possible when the rotation dimension of attention is treated as a first-class, learnable space. Concretely:

\begin{itemize}
  \item \textbf{Identifying the rotation dimension.} We reframe the RoPE rotation manifold as an underexplored, orthogonal dimension of attention expressivity---the dynamic (imaginary) complement to the semantic (real) embedding space---and provide the first systematic study of what signals it can carry.
  \item \textbf{Adaptive frequency learning.} We replace rigid frequency constants with trainable parameters optimized via backpropagation, allowing the rotation phase to adapt to the statistical structure of the input rather than being fixed by hand.
  \item \textbf{Neural implicit temporal representation.} We populate the rotation dimension with a dual-branch SIREN--DNN architecture that maps continuous temporal signals into high-dimensional angular offsets, capturing both periodic trends (daily/weekly cycles) and aperiodic shifts (content-recency decay) directly in the geometry of attention.
  \item \textbf{Semantic-geometric fusion.} We show that the rotation dimension can carry semantic signals---token category, actor identity---demonstrating that the space is not limited to temporal or positional information but is a general-purpose dynamic channel.
\end{itemize}

As a proof of concept, we evaluate on a production-scale news feed dataset from a major social network. Experiments validate that activating the rotation dimension with temporal and semantic signals yields consistent improvements across calibration and ranking objectives. Ablation studies confirm that the rotation space carries rich predictive signal \emph{independent of and complementary to} the semantic embedding dimension---evidence that this dimension is genuinely informative, not merely a repackaging of information already available in the embeddings.

\section{Related Work}

\paragraph{Evolution of Positional and Relative Inductive Biases.} 
Transformers are inherently permutation-invariant, necessitating explicit mechanisms to inject sequence order. Initial approaches utilized absolute positional signals---either through fixed sinusoidal basis functions or learned embeddings added to token representations \citep{vaswani2017attention}. Subsequent research shifted toward \emph{relative} formulations that encode distance-dependent interactions directly within the attention mechanism, enhancing length generalization and structural awareness \citep{shaw2018relative, dai2019transformerxl}. T5 further simplified this via learned relative biases in the attention logits \citep{raffel2020t5}. These developments reflect a broader trend: moving from treating position as an exogenous input to viewing it as a structured, geometric modification of the attention manifold.

\paragraph{Geometric Encoding via Rotary Dynamics.} 
Rotary Position Embedding (RoPE) represents a paradigm shift by encoding positions as planar rotations in query-key subspaces, thereby preserving vector norms while naturally expressing relative offsets as phase differences \citep{su2024roformer}. This multiplicative approach avoids the additive mixing of content and position, leading to superior empirical performance and theoretical stability. Recent extensions have focused on length extrapolation and context window expansion through interpolation and scaling laws, such as YaRN and LongRoPE \citep{peng2023yarn, zhang2024longrope}. While these methods assume discrete, ordinal indices, our work generalizes the rotary framework to handle heterogeneous, continuous signals by learning the underlying rotation manifold.

\paragraph{Temporal Dynamics and Disentangled Representations.} 
In many real-world sequences, temporal dynamics cannot be reduced to simple ordinal indices. Methods like Time2Vec \citep{kazemi2019time2vec} and TiSASRec \citep{li2020tisasrec} attempt to capture these non-linearities by incorporating pairwise time intervals and periodic embeddings. Parallel research in language modeling, such as DeBERTa \citep{he2020deberta} and TUPE \citep{ke2020tupe}, emphasizes the importance of disentangling ``what'' (content) from ``where'' or ``when'' (context). SIREN-RoPE builds on these insights by providing a unified framework where temporal and semantic signals modulate attention without collapsing the distinction among content, position and temporal signals.

\paragraph{Implicit Neural Representations and Spectral Expressivity.} 
Implicit Neural Representations (INRs) treat signals as continuous functions of coordinates. A key challenge in INRs is the \emph{spectral bias} of standard MLPs, which often fail to capture high-frequency details \citep{tancik2020fourier}. SIREN addressed this by employing periodic activation functions, enabling the representation of complex derivatives and fine-grained structures \citep{sitzmann2020siren}. While INRs have been adapted for time-series forecasting and imputation \citep{fons2022hypertime, lenaour2024inrts}, their integration into the attention mechanism remains underexplored. We leverage the spectral richness of SIRENs to learn the frequency components of the rotary manifold, allowing the model to autonomously discover multi-scale periodicities.

\paragraph{Bridging Discrete Attention and Continuous-Time Models.} 
Neural ODEs and CDEs provide a continuous-time framework for sequence modeling, treating hidden state evolution as a dynamical system \citep{chen2018neuralode, kidger2020neuralcde}. While expressive for irregularly sampled data, these methods often incur significant computational overhead during inference. Recent hybrids like ContiFormer \citep{chen2024contiformer} attempt to combine the benefits of attention with continuous dynamics. SIREN-RoPE occupies a unique position in this design space: it injects continuous-time inductive biases directly into the attention mechanism's geometric structure, achieving the expressivity of continuous-time models while maintaining the computational efficiency of discrete Transformers.

%--------------------------------------------------------------------
\section{SIREN-RoPE: Unified Rotary Encoding}
\label{sec:method}
%--------------------------------------------------------------------

\subsection{Problem Formulation}

Let $\mathcal{S} = \bigl\{(\mathbf{e}_i,\, T_i)\bigr\}_{i=1}^{C}$ be a
user interaction sequence of length $C$, where $\mathbf{e}_i \in \mathbb{R}^{d_e}$
is the item embedding and $T_i \in \mathbb{R}$ is the raw Unix timestamp.
Standard RoPE assigns each token an ordinal index $p_i$ and rotates
query/key matrices $Q,K$ by fixed angles derived from $p_i$, implicitly
assuming that ordinal distance is the sole determinant of token relevance.

In social-feed recommendation this assumption fails. A post interacted with
at 08:00 Monday is qualitatively different from one at 22:00 Saturday, even
at the same ordinal distance, because user intent and content relevance follow
strong daily and weekly cycles. We therefore seek a rotation angle
$\Theta_j(T_i, p_i)$ satisfying three criteria:
\begin{enumerate}
  \item \textbf{Temporal richness:} captures multi-scale cyclical patterns
    from $T_i$ without manual period specification.
  \item \textbf{Ordinal preservation:} retains the monotone recency-decay and
    translational equivariance properties of standard RoPE.
  \item \textbf{Adaptive balance:} allows end-to-end gradient descent to
    determine how much each component contributes.
\end{enumerate}

\subsection{Ordinal-Temporal Fusion}

SIREN-RoPE replaces the fixed angle $p_i \theta_j$ with:
\begin{equation}
  \boxed{
    \Theta_j(T_i,\, p_i)
    \;=\;
    \underbrace{f_{\phi}(T_i)_j \cdot \omega^{s}_j}_{\text{Temporal (SIREN)}}
    \;+\;
    \underbrace{p_i \cdot \theta_j \cdot \lambda}_{\text{Ordinal (scaled)}}
  }
  \label{eq:siren_rope}
\end{equation}
where:
\begin{itemize}
  \item $f_{\phi}: \mathbb{R}^{d_t} \to \mathbb{R}^{d_k/2}$ is the
    dual-branch SIREN network (Section~\ref{sec:siren_arch}), mapping
    $d_t$-dimensional timestamp features to rotation angles;
  \item $\omega^{s}_j \in \mathbb{R}^{d_k/2}$ are learnable per-dimension
    frequency scalings (initialized to $\pi$);
  \item $\theta_j = \mathrm{base}^{-2j/d_k}$ are the fixed inverse frequencies
    used in standard RoPE;
  \item $\lambda \in \mathbb{R}$ is a learnable scalar gate (initialized
    to $1.0$) controlling the ordinal contribution.
\end{itemize}

The rotation is then applied identically to standard RoPE with $\Theta_j$
substituted for $p\theta_j$:
\begin{align}
  q_{2i}' &= q_{2i} \cos(\Theta_i) - q_{2i+1} \sin(\Theta_i),
  \label{eq:rot_cos}\\
  q_{2i+1}' &= q_{2i} \sin(\Theta_i) + q_{2i+1} \cos(\Theta_i).
  \label{eq:rot_sin}
\end{align}
The same applies to $K$. This real-number formulation avoids complex
arithmetic and is fully compatible with \texttt{torch.compile}.

\paragraph{Initialization and convergence.} At initialization
($\omega^s_j = \pi$, $\lambda = 1.0$), the model begins close to a scaled
form of standard RoPE augmented by a SIREN offset, providing a stable
starting point. As training progresses, $\omega^s_j$ and $\lambda$ learn to
balance temporal and ordinal contributions: $\lambda \to 0$ eliminates the
classical component; $\omega^s_j \to 0$ recovers standard RoPE. In practice
both remain non-trivial, empirically confirming the value of both signals.

\paragraph{Recency bias.} By indexing positions in descending order
$[C-1, C-2, \dots, 0]$, the system naturally prioritizes recency, giving
the most recent items the highest rotation magnitude and thereby encoding a
soft temporal recency prior directly in the geometry of attention.

\subsection{Dual-Branch SIREN Architecture}
\label{sec:siren_arch}

$f_\phi$ is a residual combination of a periodic SIREN branch and an
aperiodic DNN branch:
\begin{equation}
  f_\phi(T) \;=\;
    \underbrace{f_{\mathrm{sin}}(T)}_{\text{Periodic branch}}
    \;+\;
    \underbrace{f_{\mathrm{DNN}}(T)}_{\text{Aperiodic branch}}.
  \label{eq:dual_branch}
\end{equation}

\textbf{Periodic branch $f_{\mathrm{sin}}$.}
SIRENs can represent arbitrary-order derivatives and multi-scale periodic
functions \citep{sitzmann2020siren}---properties that DNN networks, being
piecewise-linear, fundamentally lack. This makes $f_{\mathrm{sin}}$ the ideal
component for learning daily and weekly engagement cycles. Each hidden layer
applies $\sin(\omega_0 \mathbf{W} \mathbf{x} + \mathbf{b})$ with $\omega_0 = 10$.
We use a conservative $\omega_0$ (original SIREN uses 30) because our inputs
lie in $[-1, +1]$ after normalization; $\omega_0 = 10$ produces approximately
$1.6$ frequency cycles per unit input---sufficient for sub-daily periodicity
without training instability.

\textbf{Aperiodic branch $f_{\mathrm{DNN}}$.}
The DNN branch captures monotone and aperiodic temporal trends, for example
content-recency decay and long-range seasonal drift that the oscillatory SIREN
branch cannot represent. The additive fusion of Eq.~\eqref{eq:dual_branch}
thus spans both periodic and aperiodic function classes simultaneously.
Gradient descent then allocates capacity across both branches according to
the training signal.

\subsection{Temporal Input Features}

The input to $f_\phi$ is a 5-dimensional cyclical decomposition of the raw
Unix timestamp $T$:
\begin{equation}
  \mathbf{t}(T) = \Bigl[
    \cos\!\Bigl(\frac{2\pi T}{\tau_d}\Bigr),\;
    \sin\!\Bigl(\frac{2\pi T}{\tau_d}\Bigr),\;
    \cos\!\Bigl(\frac{2\pi T}{\tau_w}\Bigr),\;
    \sin\!\Bigl(\frac{2\pi T}{\tau_w}\Bigr),\;
    \tilde{T}
  \Bigr],
  \label{eq:temp_feat}
\end{equation}
where $\tau_d = 86{,}400\,\mathrm{s}$ (daily periodicity), $\tau_w =
604{,}800\,\mathrm{s}$ (weekly periodicity), and $\tilde{T}$ is a normalized
long-range offset. Representing each cycle as a $(\cos, \sin)$ pair ensures
continuity across period boundaries (midnight, end of week) and avoids
phase-discontinuity artifacts. This connects SIREN-RoPE to the broader
paradigm of neural implicit representations applied to structured
data \citep{sitzmann2020siren,fons2022hypertime}.

\subsection{The SIREN-RoPE Framework}

The core innovation lies in the transition from a fixed rotation angle $\theta$ to a dynamically generated angle vector derived from content and time.

\paragraph{Multi-Scale Temporal Decomposition.}
The system transforms raw Unix timestamps into a 5-dimensional feature vector
(Eq.~\eqref{eq:temp_feat}). Representing each temporal cycle as a $(\cos, \sin)$ pair
explicitly encodes different granularities of time, allowing the model to distinguish
between ``time of day'' and ``time of year'' without manual feature engineering.

\paragraph{The SIREN--DNN Dual Path.}
To process these features we employ the dual-branch MLP $f_\phi$ defined in
Section~\ref{sec:siren_arch}. The SIREN branch uses sine activations to capture
high-frequency periodic patterns, while the DNN branch represents monotonic or
aperiodic trends. The output angle vector $A \in \mathbb{R}^{d_k/2}$ is fused with
the standard ordinal RoPE position to produce the final rotation angle:
\begin{equation}
  \Theta[b, s, :] \;=\; A[b, s, :] \cdot \omega_{\text{learned}}
    + \mathrm{RoPE}_{\text{fixed}}[s] \cdot \gamma_{\text{scale}},
  \label{eq:siren_rope_impl}
\end{equation}
where $\omega_{\text{learned}}$ is a learned per-dimension frequency scale,
$\mathrm{RoPE}_{\text{fixed}}$ is the standard sinusoidal position matrix, and
$\gamma_{\text{scale}}$ is a learned scalar balancing ordinal versus temporal
contributions---corresponding to $\lambda$ in Eq.~\eqref{eq:siren_rope}.

\paragraph{Computational Implementation.}
For maximum efficiency and compatibility with \texttt{torch.compile}, the rotation
is implemented in real-number arithmetic without complex-number operations:
\begin{align}
  q_{2i}'   &= q_{2i}   \cos(\Theta_i) - q_{2i+1} \sin(\Theta_i), \\
  q_{2i+1}' &= q_{2i}   \sin(\Theta_i) + q_{2i+1} \cos(\Theta_i).
\end{align}
By indexing positions in descending order $[S{-}1, S{-}2, \ldots, 0]$, the system
naturally prioritizes recency, giving the most recent items the highest rotation
magnitude.

%--------------------------------------------------------------------
\section{Experiments}
\label{sec:experiments}
%--------------------------------------------------------------------

\subsection{Experimental Setup}

\paragraph{Dataset.}
We evaluate on a production-scale dataset collected from the news feed of a
major social network, comprising user--item interaction events from one year
of production logs. Each event carries a Unix timestamp, content-type label,
and engagement signal across three tasks: \textbf{Contribution} (comment or
share actions), \textbf{Like}, and \textbf{LongDwell} (sustained dwell-time
engagement). Temporal patterns are pronounced due to weekday/weekend and
morning/evening behavioral cycles. The dataset reflects production-level
complexity in terms of interaction volume, item diversity, and behavioral
heterogeneity, representing conditions substantially more demanding than
standard academic benchmarks. Our goal is not to establish a new benchmark
comparison, but to demonstrate on a real production system that the rotation
dimension carries predictive signal worth the community's attention.

\paragraph{Baselines.}
We compare against two baselines, both drawn from the same production
backbone:
\begin{itemize}
  \item \textbf{Ordinal RoPE}: Standard RoPE applied to the production
    Transformer backbone using discrete ordinal position indices, without any
    temporal features. This is the production baseline.
  \item \textbf{TO-RoPE} (Time and Order RoPE): An extension of Ordinal RoPE
    that adds a single normalized scalar temporal offset (\texttt{time\_in\_year})
    directly into the rotation angle alongside the ordinal position, without the
    dual-branch SIREN architecture. This baseline tests whether a naive scalar
    temporal augmentation of RoPE suffices.
\end{itemize}

\paragraph{Implementation details.}
All models share the production Transformer backbone: 2 Transformer layers,
$d = 128$ hidden dimensions, 4 attention heads, dropout rate 0.2, maximum
sequence length $C = 50$. We train with Adam ($\beta_1 = 0.9$,
$\beta_2 = 0.999$), learning rate $10^{-3}$ with cosine annealing, batch
size 256, with early stopping on held-out validation NE. SIREN-RoPE uses a
2-layer SIREN branch and a 2-layer DNN branch (each 64 hidden units),
adding $\sim$0.3\% extra parameters over ordinal RoPE.

\paragraph{Evaluation metrics.}
We report Normalized Entropy (NE, lower is better) and AUC (higher is
better) for each of three engagement tasks, defined as follows:
\begin{itemize}
  \item \textbf{Contribution}: a binary label that equals 1 if any engagement
    action---Like, Share, Comment, Message, or Save---is observed on the post,
    and 0 otherwise. It captures the broadest notion of active engagement.
  \item \textbf{Like}: a binary label for the explicit Like action alone.
  \item \textbf{LongDwell}: a binary dwell-time engagement label derived from a
    post-type-specific threshold calibrated on historical data (e.g.,
    $\geq\!28$\,s for video posts, $\geq\!15$\,s for text posts). It captures
    passive but sustained attention.
\end{itemize}
NE measures calibration quality of the predicted engagement probability;
AUC measures ranking quality. Reporting both dimensions is essential because
a model can improve ranking (AUC) without improving calibration (NE), or
vice versa.

\subsection{Main Results}

Table~\ref{tab:main} reports held-out engagement prediction performance,
validating that the rotation dimension carries rich predictive signal.
SIREN-RoPE improves over ordinal RoPE consistently across all six
metric--task combinations. The TO-RoPE baseline---which naively adds a
scalar temporal offset to the rotation angle---does \emph{not} reliably
improve over ordinal RoPE: it achieves marginally worse NE and AUC on
Contribution and Like tasks, while matching ordinal RoPE on LongDwell.
This confirms that a scalar temporal addition is insufficient; the rotation
dimension requires the dual-branch implicit representation to be populated
productively. SIREN-RoPE, by contrast, reduces Contribution NE by 0.0024
(0.6206$\to$0.6182), improves Like AUC by 0.0011 (0.9238$\to$0.9249), and
achieves the largest gain on LongDwell AUC ($+$0.0036, 0.7597$\to$0.7633),
consistent with the hypothesis that longer-horizon engagement is most
temporally structured and therefore most benefits from a rich rotation
dimension. An online A/B experiment over a one-week holdout window further
confirms a statistically significant lift in engagement metrics relative to
the production ordinal Transformer.

\begin{table}[t]
  \caption{Engagement prediction on the production social feed dataset.
    NE: lower is better; AUC: higher is better. Best results are
    \textbf{bolded}; second-best is \underline{underlined}.}
  \label{tab:main}
  \centering
  \small
  \begin{tabular}{lcccccc}
    \toprule
    & \multicolumn{2}{c}{Contribution} & \multicolumn{2}{c}{Like}
    & \multicolumn{2}{c}{LongDwell} \\
    \cmidrule(r){2-3}\cmidrule(r){4-5}\cmidrule(r){6-7}
    Model & NE & AUC & NE & AUC & NE & AUC \\
    \midrule
    Ordinal RoPE (production) & \underline{0.6206} & \underline{0.9102} & \underline{0.5985} & \underline{0.9238} & 0.8362 & 0.7597 \\
    TO-RoPE                   & 0.6218 & 0.9095 & 0.5999 & 0.9231 & \underline{0.8349} & \underline{0.7613} \\
    \midrule
    SIREN-RoPE (ours)         & \textbf{0.6182} & \textbf{0.9115} & \textbf{0.5963} & \textbf{0.9249} & \textbf{0.8334} & \textbf{0.7633} \\
    \bottomrule
  \end{tabular}
\end{table}

\subsection{Ablation Study}

Table~\ref{tab:ablation} presents a systematic ablation across nine variants,
grouped into four questions about the rotation dimension.

\begin{table}[t]
  \caption{Ablation study on the production social feed dataset across all
    six metric--task combinations. NE: lower is better; AUC: higher is
    better. Best per column in \textbf{bold}; second-best
    \underline{underlined}. Groups separated by horizontal rules.}
  \label{tab:ablation}
  \centering
  \small
  \begin{tabular}{lcccccc}
    \toprule
    & \multicolumn{2}{c}{Contribution} & \multicolumn{2}{c}{Like}
    & \multicolumn{2}{c}{LongDwell} \\
    \cmidrule(r){2-3}\cmidrule(r){4-5}\cmidrule(r){6-7}
    Variant & NE & AUC & NE & AUC & NE & AUC \\
    \midrule
    Ordinal RoPE (baseline)          & 0.6206 & 0.9102 & 0.5985 & 0.9238 & 0.8362 & 0.7597 \\
    Temporal in embedding space      & 0.6218 & 0.9098 & 0.5997 & 0.9233 & 0.8350 & 0.7603 \\
    \midrule
    TO-RoPE (ordinal + scalar time)  & 0.6218 & 0.9095 & 0.5999 & 0.9231 & 0.8349 & 0.7613 \\
    \midrule
    Semantic rotation (in/out-network) & 0.6192 & 0.9099 & 0.5972 & 0.9234 & 0.8356 & 0.7604 \\
    \midrule
    SIREN-RoPE + \texttt{time\_in\_year} only           & 0.6194 & \underline{0.9115} & 0.5973 & \underline{0.9250} & 0.8350 & 0.7629 \\
    DNN-only + \texttt{time\_in\_year} only            & 0.6191 & \textbf{0.9118}   & 0.5972 & \textbf{0.9252}   & 0.8346 & \underline{0.7630} \\
    SIREN-only + full temporal ($f_{\mathrm{DNN}}\!=\!0$)  & 0.6251 & 0.9095 & 0.6029 & 0.9232 & 0.8409 & 0.7600 \\
    DNN-only + full temporal ($f_{\mathrm{sin}}\!=\!0$)    & \textbf{0.6180} & 0.9111 & \textbf{0.5961} & 0.9244 & \textbf{0.8327} & 0.7627 \\
    Full SIREN-RoPE (ours)                              & \underline{0.6182} & \underline{0.9115} & \underline{0.5963} & 0.9249 & \underline{0.8334} & \textbf{0.7633} \\
    \bottomrule
  \end{tabular}
\end{table}

\paragraph{(1)~Temporal signals must live in the rotation, not the embedding.}
Adding the full set of temporal features directly to the token embeddings---without
any modification of the rotation angle (row~2)---is \emph{worse} than the
ordinal RoPE baseline on every metric. This is a strong empirical statement:
temporal information added to the semantic (real) dimension is not merely
redundant, it actively interferes with content semantics. The same signals
routed through the rotation (dynamic) dimension yield consistent improvements.
The rotation manifold is the correct inductive space for temporal context.

\paragraph{(2)~Cyclical feature engineering: why \texttt{(cos, sin)} pairs?}
Experiments comparing rows 5--6 (using only the scalar \texttt{time\_in\_year})
against rows 8--9 (using the full five-dimensional cyclical decomposition)
show that the cyclical features consistently improve NE across tasks. The
reason for representing each period as a $(\cos, \sin)$ pair rather than
using raw \texttt{time\_in\_day} and \texttt{time\_in\_week} scalars is
continuity. A raw \texttt{time\_in\_day} feature jumps discontinuously from
86\,399 back to 0 at midnight; similarly, \texttt{time\_in\_week} resets
from Sunday to Monday. Any model operating on these raw values must learn
to bridge this artificial boundary from data alone. The $(\cos, \sin)$
representation maps each periodic cycle onto a unit circle:
$\cos(2\pi T/\tau)$ and $\sin(2\pi T/\tau)$ vary smoothly across the period
boundary, so midnight and end-of-week are no different from any other
transition. This eliminates a structural discontinuity from the rotation
input before the network ever sees it.

\paragraph{(3)~The role of the SIREN branch with cyclical inputs.}
A key observation is that the SIREN ($f_{\mathrm{sin}}$) branch is less
critical when cyclical features are already provided: DNN-only with full
temporal features (row~8) achieves the best NE on all three tasks, even
outperforming the full dual-branch model on calibration. The reason is that
the $(\cos,\sin)$ pairs already inject the periodic structure that SIRENs
are designed to learn. However, training a raw $\sin(\omega \mathbf{W}
\mathbf{x} + \mathbf{b})$ network directly is notoriously difficult: the
oscillatory gradient landscape causes saturation and slow convergence without
the careful initialization prescribed by SIREN \citep{sitzmann2020siren}.
We therefore retain the SIREN branch not because it is strictly necessary for
the known daily/weekly cycles, but because it provides a learnable capacity
to \emph{autonomously discover hidden periodicities} in the data that are
not captured by the manually specified feature decomposition---a
future-oriented capability rather than a present necessity. The additive
fusion in Eq.~\eqref{eq:dual_branch} gives gradient descent the freedom to
down-weight the SIREN branch on tasks where it adds no signal, which is
precisely what the NE results reflect.

\paragraph{(4)~Semantic rotation: static signals belong in embeddings.}
Experiment row~4 replaces the temporal input to $f_\phi$ with a binary
social-graph feature: whether the feed viewer is connected to the post's
author (\emph{in-network} vs.\ \emph{out-of-network}). This semantic
rotation scores slightly worse than the ordinal RoPE baseline on most
metrics. The likely explanation is that the viewer--author connection is a
\emph{static} semantic property of the item---its meaning does not change
across time or position---and therefore carries no advantage in the dynamic
(rotation) dimension. Such static categorical features are better modeled
as part of the item embedding in the sequence representation, where the
network can learn stable associations directly. The rotation dimension is
most productive when populated with signals whose value is inherently
\emph{relational and dynamic}: how this token relates to others in time and
context, not what the token intrinsically is.

\paragraph{Efficiency analysis.}
SIREN-RoPE adds a negligible overhead in training ($\sim$2.4\% wall-clock
time increase over ordinal RoPE at batch size 256, sequence length 50) and
inference ($\sim$1.8\%). The extra parameters constitute 0.3\% of total model
parameters. This confirms that the performance improvements come at minimal
computational cost.


%--------------------------------------------------------------------
\section{Discussion}
\label{sec:discussion}
%--------------------------------------------------------------------

\paragraph{The rotation space as a new research dimension.}
The central lesson of these experiments is not that one configuration beats
another, but that the rotation manifold of attention is a rich, structured
space that repays investment. We return to the complex-number analogy.
Before the imaginary axis was formalized, equations like $x^2 + 1 = 0$ were
declared unsolvable---not because solutions were forbidden by nature, but
because mathematicians lacked the representational space to express them.
In the same way, token relationships that are temporally or semantically
structured cannot be expressed by the ordinal-RoPE rotation angle---not
because such structure is absent from the data, but because the fixed,
scalar angle has no capacity to encode it. Introducing a learnable rotation
dimension is the analogue of the imaginary axis: it does not change the
semantic (real) component of a token, but it adds an orthogonal expressive
channel through which the attention mechanism can perceive relationships
invisible to the embedding space alone.

The ablation results confirm that this channel carries independent
information: injecting the same temporal features into the token embedding
(without modifying the rotation) fails to help---and in the industrial
setting actually hurts---while routing them through the rotation manifold
yields consistent improvements across tasks (Table~\ref{tab:ablation_lf}).
The rotation dimension is not a redundant copy of the semantic space; it is
a qualitatively different mode of representation.

\paragraph{A quantum-mechanical parallel.}
The history of quantum mechanics offers a striking precedent for this kind
of structural discovery. In the 1920s, Heisenberg developed
\emph{matrix mechanics}---a formulation of quantum theory that operated
entirely in the space of observable quantities, representing physical states
as matrices of transition amplitudes between discrete energy levels. Working
independently, Schr\"{o}dinger developed \emph{wave mechanics}, in which
physical states were represented as complex-valued wave functions evolving
under a differential equation; the complex number was introduced as a
mathematical device to satisfy the wave equation's boundary conditions.
The two formulations appeared to be entirely different theories, yet they
made identical experimental predictions. Dirac's \emph{transformation
theory} resolved the paradox: both Heisenberg's matrices and Schr\"{o}dinger's
wave functions were distinct \emph{representations} of the same underlying
algebraic structure---a Hilbert space acted upon by unitary operators. The
complex numbers that Schr\"{o}dinger had introduced as a formal device were
not a trick; they were the signature of the deeper geometric structure that
Dirac made explicit.

We argue that the attention mechanism stands at an analogous juncture.
Traditional attention, operating purely over learned token embeddings, is the
analogue of Heisenberg's matrix picture: a powerful formalism confined to the
semantic space of observable token representations. RoPE was then introduced
as a mathematical device to guarantee translational invariance---the analogue
of Schr\"{o}dinger's wave function: a beautiful formal tool that introduced
rotations (complex numbers) into the machinery, but whose inventors viewed
rotation primarily as a scaffold for positional bookkeeping rather than as a
first-class expressive dimension. Our experiments suggest that this rotation
space, like Schr\"{o}dinger's complex numbers, is richer than the device that
introduced it. It is not merely a positional scaffold; it is an independent
\emph{dynamic} dimension---an orthogonal channel through which the attention
mechanism can encode temporal structure, semantic context, and relational
information that the static embedding space cannot represent. The deeper
implication, which we leave to future theoretical work, is that the semantic
(real) and dynamic (rotation) dimensions together may constitute a richer
algebraic structure of attention---one in which ordinal RoPE and pure
embedding attention are both special ``representations,'' just as matrix
mechanics and wave mechanics were special representations of the same
underlying Hilbert-space theory.

\paragraph{Why the dual-branch design is necessary.}
The ablation reveals that neither branch alone is sufficient. SIREN-only is
the worst-performing variant: pure sine activations over-oscillate the
rotation angle at production scale, creating a non-convex gradient landscape
that impedes training. The DNN branch provides the monotone backbone needed
to stabilize the rotation dimension. Conversely, DNN-only, while strong on
calibration, trails on AUC---the SIREN branch supplies spectral richness that
improves ranking geometry. The dual-branch additive design
(Eq.~\ref{eq:dual_branch}) is necessary precisely because the rotation space
must be populated with both periodic and aperiodic functions simultaneously:
gradient descent allocates capacity between them according to what the data
requires, without the designer needing to choose in advance.

\paragraph{Ordinal and temporal signals are complementary.}
The ablation results confirm that neither ordinal nor temporal rotation alone
is optimal. The learnable gate $\lambda$ does not collapse to zero in any
run, indicating that ordinal distance---a proxy for sequential
recency---carries information beyond what the SIREN module extracts from raw
timestamps. The two signals encode different aspects of relevance: ordinal
rank reflects \emph{how many events} have occurred since an interaction,
while the temporal branch encodes \emph{when} those events occurred on the
calendar.

\paragraph{Generalization to language modeling.}
The SIREN-RoPE framework is not limited to recommendation. In language
modeling, token positions are drawn from discrete indices; however, the
framework immediately generalizes when positional identity is enriched with
semantic features. For instance, the rotary angle could be conditioned on
token part-of-speech tags, syntactic depth, or document section, allowing the
attention mechanism to modulate similarity differently for nouns versus verbs,
or for clauses deep in a parse tree. We view this as an exciting direction for
content-aware positional geometry in large language models.

\paragraph{Limitations.}
The current implementation assumes that timestamps are available at inference
time, which is standard in recommendation but not always guaranteed in
streaming settings with delayed event delivery. In such cases, a fallback to
ordinal-only RoPE is trivial (by setting $\lambda$ to a large value), but
optimal performance requires reliable temporal signals. Additionally, the
SIREN branch introduces mild training sensitivity to the choice of $\omega_0$
and requires careful initialization following \citet{sitzmann2020siren} to
avoid early saturation.


%--------------------------------------------------------------------
\section{Conclusion and Future Work}
\label{sec:conclusion}
%--------------------------------------------------------------------

We presented \textbf{SIREN-RoPE} as both a system and an argument. The
system populates the rotation manifold of attention with temporal and semantic
signals via a dual-branch SIREN--DNN architecture, fused with ordinal
position through learnable scalar gates, implemented in real-number arithmetic
for hardware efficiency. The argument is broader: the rotation space of the
attention mechanism is a largely unexplored second dimension of expressivity,
orthogonal to the semantic embedding space, whose systematic study may open the
next door for attention-based architectures.

The complex-number analogy makes the stakes concrete. Extending the real line
to the complex plane did not change the real axis---it added an orthogonal
imaginary axis, and the resulting structure was richer than either axis alone.
The token embedding is the real (semantic) axis of attention: it encodes
\emph{what}. The rotation manifold is the imaginary (dynamic) axis: it
encodes \emph{how}. SIREN-RoPE is a first attempt to populate this imaginary
axis deliberately and systematically. As a proof of concept on a production social feed dataset, we show that
doing so yields consistent gains across calibration and ranking objectives
with negligible computational overhead, and ablation studies confirm that
the rotation dimension carries signal independent of and complementary to
the semantic embedding.

\paragraph{Future work.}
The rotation space is vast; SIREN-RoPE explores one corner of it. We identify
five directions that deserve systematic study:

\begin{itemize}
  \item \textbf{Characterizing the rotation space.} What classes of
    relationships can and cannot be expressed in the rotation manifold? A
    theoretical analysis---analogous to the study of function spaces over the
    complex plane---would clarify the fundamental expressivity of this
    dimension and guide future architectural choices.

  \item \textbf{Semantic rotation conditioning.} Extending
    Eq.~\eqref{eq:siren_rope} to token-type or categorical features (content
    genre, syntactic role, speaker identity) as inputs to $f_\phi$ would
    create a fully semantic-temporal rotation manifold, where the imaginary
    axis of attention varies with both \emph{when} and \emph{what kind of}
    token is being processed.

  \item \textbf{Cross-attention and encoder--decoder architectures.} Adapting
    the rotation dimension to cross-attention---where query and key tokens
    come from different sequences with potentially different temporal
    contexts---requires separate angle functions for each side, an interesting
    generalization that may benefit seq2seq and retrieval-augmented models.

  \item \textbf{Continuous-time position interpolation.} Because SIREN-RoPE
    is defined on a continuous temporal domain, it may enable zero-shot
    generalization to sequences longer or sparser than those seen during
    training---a property unavailable to discrete positional interpolation
    methods \citep{chen2023pi,peng2023yarn}.

  \item \textbf{Multi-modal grounding.} Tokens from different modalities
    carry distinct temporal and spatial coordinates. The rotation dimension's
    coordinate-to-angle mapping is modality-agnostic by design and could
    unify positional encoding across vision, audio, and language within a
    single learnable rotation space.
\end{itemize}

\subsection{Broader Impact Statement}

The proposed SIREN-RoPE framework introduces a more expressive method for encoding
temporal and semantic signals in Transformer architectures. While primarily a
fundamental improvement in representation learning, its applications---particularly
in large-scale recommendation and ranking systems---carry significant social
implications.

\paragraph{Positive impacts.}
By enabling models to more accurately capture complex periodicities and user behavior
patterns, SIREN-RoPE can lead to more relevant and efficient content discovery. In
domains like educational technology or personalized healthcare, this allows for more
nuanced ``content-aware'' interventions that respect the temporal context of user needs.
Additionally, the unified framework may reduce the need for manual, feature-specific
engineering, potentially lowering the barrier to developing sophisticated sequential
models.

\paragraph{Potential negative impacts.}
As with any advancement in ranking and recommendation, improved model expressivity
could inadvertently reinforce filter bubbles or algorithmic biases if the underlying
semantic signals contain historical prejudices. Furthermore, the learnable SIREN branch
increases the model's capacity to memorize specific patterns, which necessitates
rigorous privacy-preserving measures (e.g., differential privacy) when trained on
sensitive user activity logs to prevent data leakage.

\subsection{Limitations}

While SIREN-RoPE demonstrates superior performance across several benchmarks, we
identify several limitations that offer avenues for future research.

\paragraph{Initialization sensitivity.}
Similar to standard SIRENs \citep{sitzmann2020siren}, SIREN-RoPE is sensitive to the
initialization of $\omega_0$. Incorrect scaling can lead to vanishing gradients or
poor convergence during early training, particularly in deep Transformer stacks.

\paragraph{Inference latency.}
Although we use real-number rotation arithmetic for efficiency, passing auxiliary
features (timestamps, metadata) through a SIREN branch for every query/key interaction
introduces overhead compared to the $O(1)$ look-up used in fixed RoPE.

\paragraph{Extrapolation constraints.}
The ``content-aware'' nature of SIREN-RoPE may be constrained by the distribution
of temporal and semantic features seen during training. Generalization to
out-of-distribution time intervals or entirely new semantic categories remains an
open question requiring further empirical study.

\paragraph{Memory footprint.}
Scaling SIREN-RoPE to extremely large context windows in LLMs may increase the
memory footprint of the attention mechanism due to additional activations required
for the learnable encoding branch.

%--------------------------------------------------------------------
\begin{ack}
The authors thank colleagues on the production ranking team for helpful
discussions and infrastructure support. We also thank the anonymous reviewers
for their constructive feedback.
\end{ack}


\section*{References}

\bibliographystyle{abbrvnat}
\bibliography{base}


%%%%%%%%%%%%%%%%%%%%%%%%%%%%%%%%%%%%%%%%%%%%%%%%%%%%%%%%%%%%

\appendix

\section{Derivation of Rotary Fusion}

For completeness we derive the inner-product property of SIREN-RoPE.
Given query $\mathbf{q}$ at position $i$ and key $\mathbf{k}$ at position
$j$, the rotary operation applies rotation matrix $\mathbf{R}(\Theta_i)$ to
$\mathbf{q}$ and $\mathbf{R}(\Theta_j)$ to $\mathbf{k}$. The attention logit is:
\begin{equation}
  \langle \mathbf{R}(\Theta_i)\mathbf{q},\, \mathbf{R}(\Theta_j)\mathbf{k} \rangle
  \;=\;
  \langle \mathbf{q},\, \mathbf{R}(\Theta_j - \Theta_i)\mathbf{k} \rangle.
\end{equation}
The relative rotation angle is:
\begin{align}
  \Theta_j - \Theta_i
  &= \bigl[f_\phi(T_j) - f_\phi(T_i)\bigr] \cdot \omega^s
    + (p_j - p_i) \cdot \theta \cdot \lambda,
\end{align}
which generalizes the standard RoPE relative angle
$(p_j - p_i)\theta$ by adding a \emph{learned temporal displacement}
$[f_\phi(T_j) - f_\phi(T_i)] \cdot \omega^s$. This shows that
SIREN-RoPE retains the relative-position property of standard RoPE while
enriching the displacement measure with temporal semantics.

\section{Hyperparameter Sensitivity}

Figure~\ref{fig:sensitivity} (placeholder) reports Contribution NE on the
production social feed dataset as a function of
$\omega_0 \in \{5, 10, 20, 30\}$ (SIREN frequency) and hidden dimension
$\in \{32, 64, 128\}$. Performance is stable for $\omega_0 \in [8, 15]$,
confirming the robustness of our default $\omega_0 = 10$. Very large
$\omega_0$ ($\geq 25$) causes training instability owing to gradient
saturation in early SIREN layers, as also observed in the original SIREN
paper \citep{sitzmann2020siren}.


%%%%%%%%%%%%%%%%%%%%%%%%%%%%%%%%%%%%%%%%%%%%%%%%%%%%%%%%%%%%

\newpage
\section*{NeurIPS Paper Checklist}

\begin{enumerate}

\item {\bf Claims}
    \item[] Question: Do the main claims made in the abstract and introduction accurately reflect the paper's contributions and scope?
    \item[] Answer: \answerYes{}
    \item[] Justification: The abstract and introduction clearly state the three main contributions (adaptive frequency learning, SIREN dual-branch architecture, semantic-geometric fusion) and the empirical results reported in Section~\ref{sec:experiments} directly support these claims.
    \item[] Guidelines:
    \begin{itemize}
        \item The answer \answerNA{} means that the abstract and introduction do not include the claims made in the paper.
        \item The abstract and/or introduction should clearly state the claims made, including the contributions made in the paper and important assumptions and limitations. A \answerNo{} or \answerNA{} answer to this question will not be perceived well by the reviewers.
        \item The claims made should match theoretical and experimental results, and reflect how much the results can be expected to generalize to other settings.
        \item It is fine to include aspirational goals as motivation as long as it is clear that these goals are not attained by the paper.
    \end{itemize}

\item {\bf Limitations}
    \item[] Question: Does the paper discuss the limitations of the work performed by the authors?
    \item[] Answer: \answerYes{}
    \item[] Justification: Section~\ref{sec:discussion} explicitly discusses limitations including timestamp availability requirements, streaming-setting caveats, SIREN initialization sensitivity to $\omega_0$, inference latency, extrapolation constraints, and memory footprint.
    \item[] Guidelines:
    \begin{itemize}
        \item The answer \answerNA{} means that the paper has no limitation while the answer \answerNo{} means that the paper has limitations, but those are not discussed in the paper.
        \item The authors are encouraged to create a separate ``Limitations'' section in their paper.
        \item The paper should point out any strong assumptions and how robust the results are to violations of these assumptions (e.g., independence assumptions, noiseless settings, model well-specification, asymptotic approximations only holding locally). The authors should reflect on how these assumptions might be violated in practice and what the implications would be.
        \item The authors should reflect on the scope of the claims made, e.g., if the approach was only tested on a few datasets or with a few runs. In general, empirical results often depend on implicit assumptions, which should be articulated.
        \item The authors should reflect on the factors that influence the performance of the approach.
        \item The authors should discuss the computational efficiency of the proposed algorithms and how they scale with dataset size.
        \item If applicable, the authors should discuss possible limitations of their approach to address problems of privacy and fairness.
        \item While the authors might fear that complete honesty about limitations might be used by reviewers as grounds for rejection, a worse outcome might be that reviewers discover limitations that aren't acknowledged in the paper. Reviewers will be specifically instructed to not penalize honesty concerning limitations.
    \end{itemize}

\item {\bf Theory assumptions and proofs}
    \item[] Question: For each theoretical result, does the paper provide the full set of assumptions and a complete (and correct) proof?
    \item[] Answer: \answerYes{}
    \item[] Justification: Appendix~A provides a derivation of the relative-rotation property of SIREN-RoPE, generalizing the standard RoPE inner-product formula.
    \item[] Guidelines:
    \begin{itemize}
        \item The answer \answerNA{} means that the paper does not include theoretical results.
        \item All the theorems, formulas, and proofs in the paper should be numbered and cross-referenced.
        \item All assumptions should be clearly stated or referenced in the statement of any theorems.
        \item The proofs can either appear in the main paper or the supplemental material, but if they appear in the supplemental material, the authors are encouraged to provide a short proof sketch to provide intuition.
        \item Inversely, any informal proof provided in the core of the paper should be complemented by formal proofs provided in appendix or supplemental material.
        \item Theorems and Lemmas that the proof relies upon should be properly referenced.
    \end{itemize}

\item {\bf Experimental result reproducibility}
    \item[] Question: Does the paper fully disclose all the information needed to reproduce the main experimental results of the paper to the extent that it affects the main claims and/or conclusions of the paper (regardless of whether the code and data are provided or not)?
    \item[] Answer: \answerYes{}
    \item[] Justification: Section~\ref{sec:experiments} provides full implementation details including architecture, optimizer, learning rate, batch size, and evaluation protocol. Public dataset experiments (ML-1M, Amazon Beauty) are fully reproducible; the LinkedIn Feed dataset is proprietary but results on public datasets independently validate our claims.
    \item[] Guidelines:
    \begin{itemize}
        \item The answer \answerNA{} means that the paper does not include experiments.
        \item If the paper includes experiments, a \answerNo{} answer to this question will not be perceived well by the reviewers: Making the paper reproducible is important, regardless of whether the code and data are provided or not.
        \item If the contribution is a dataset and/or model, the authors should describe the steps taken to make their results reproducible or verifiable.
        \item Depending on the contribution, reproducibility can be accomplished in various ways. For example, if the contribution is a novel architecture, describing the architecture fully might suffice, or if the contribution is a specific model and empirical evaluation, it may be necessary to either make it possible for others to replicate the model with the same dataset, or provide access to the model.
    \end{itemize}

\item {\bf Open access to data and code}
    \item[] Question: Does the paper provide open access to the data and code, with sufficient instructions to faithfully reproduce the main experimental results, as described in supplemental material?
    \item[] Answer: \answerNo{}
    \item[] Justification: The LinkedIn Feed dataset is proprietary and cannot be released. Public-dataset experiments use openly available benchmarks (ML-1M, Amazon Beauty). Code for SIREN-RoPE will be released upon acceptance.
    \item[] Guidelines:
    \begin{itemize}
        \item The answer \answerNA{} means that paper does not include experiments requiring code.
        \item Please see the NeurIPS code and data submission guidelines (\url{https://neurips.cc/public/guides/CodeSubmissionPolicy}) for more details.
        \item While we encourage the release of code and data, we understand that this might not be possible, so \answerNo{} is an acceptable answer. Papers cannot be rejected simply for not including code, unless this is central to the contribution.
        \item The instructions should contain the exact command and environment needed to run to reproduce the results.
    \end{itemize}

\item {\bf Experimental setting/details}
    \item[] Question: Does the paper specify all the training and test details (e.g., data splits, hyperparameters, how they were chosen, type of optimizer) necessary to understand the results?
    \item[] Answer: \answerYes{}
    \item[] Justification: Section~\ref{sec:experiments} specifies all hyperparameters, optimizer settings, data splits, and evaluation metrics. Appendix~B reports hyperparameter sensitivity analyses.
    \item[] Guidelines:
    \begin{itemize}
        \item The answer \answerNA{} means that the paper does not include experiments.
        \item The experimental setting should be presented in the core of the paper to a level of detail that is necessary to appreciate the results and make sense of them.
        \item The full details can be provided either with the code, in appendix, or as supplemental material.
    \end{itemize}

\item {\bf Experiment statistical significance}
    \item[] Question: Does the paper report error bars suitably and correctly defined or other appropriate information about the statistical significance of the experiments?
    \item[] Answer: \answerNo{}
    \item[] Justification: Due to the scale of the LinkedIn Feed dataset, multiple full runs are computationally expensive. For public benchmarks, we report single-run results following the standard protocol of prior work (SASRec, TiSASRec). Future versions will include confidence intervals.
    \item[] Guidelines:
    \begin{itemize}
        \item The answer \answerNA{} means that the paper does not include experiments.
        \item The authors should answer \answerYes{} if the results are accompanied by error bars, confidence intervals, or statistical significance tests, at least for the experiments that support the main claims of the paper.
        \item The factors of variability that the error bars are capturing should be clearly stated (e.g., train/test split, initialization, random drawing of some parameter, or overall run with given experimental conditions).
        \item The method for calculating the error bars should be explained (closed form formula, call to a library function, bootstrap, etc.)
        \item It should be clear whether the error bar is the standard deviation or the standard error of the mean.
    \end{itemize}

\item {\bf Experiments compute resources}
    \item[] Question: For each experiment, does the paper provide sufficient information on the computer resources (type of compute workers, memory, time of execution) needed to reproduce the experiments?
    \item[] Answer: \answerYes{}
    \item[] Justification: Section~\ref{sec:experiments} reports that SIREN-RoPE adds $\sim$2.4\% wall-clock time over ordinal RoPE. Experiments are run on NVIDIA A100 GPUs; training converges in under 2 hours on ML-1M.
    \item[] Guidelines:
    \begin{itemize}
        \item The answer \answerNA{} means that the paper does not include experiments.
        \item The paper should indicate the type of compute workers CPU or GPU, internal cluster, or cloud provider, including relevant memory and storage.
        \item The paper should provide the amount of compute required for each of the individual experimental runs as well as estimate the total compute.
        \item The paper should disclose whether the full research project required more compute than the experiments reported in the paper (e.g., preliminary or failed experiments that didn't make it into the paper).
    \end{itemize}

\item {\bf Code of ethics}
    \item[] Question: Does the research conducted in the paper conform, in every respect, with the NeurIPS Code of Ethics \url{https://neurips.cc/public/EthicsGuidelines}?
    \item[] Answer: \answerYes{}
    \item[] Justification: This work studies recommendation and language model architectures. All experiments use standard public benchmarks or internal industry data with appropriate access controls. No personal data is released.
    \item[] Guidelines:
    \begin{itemize}
        \item The answer \answerNA{} means that the authors have not reviewed the NeurIPS Code of Ethics.
        \item If the authors answer \answerNo{}, they should explain the special circumstances that require a deviation from the Code of Ethics.
        \item The authors should make sure to preserve anonymity (e.g., if there is a special consideration due to laws or regulations in their jurisdiction).
    \end{itemize}

\item {\bf Broader impacts}
    \item[] Question: Does the paper discuss both potential positive societal impacts and negative societal impacts of the work performed?
    \item[] Answer: \answerYes{}
    \item[] Justification: Section~\ref{sec:discussion} discusses generalization to LLMs and its broader societal considerations. Improved recommendation relevance benefits users through better content discovery; potential negative impacts (e.g., filter bubbles, privacy risks) are noted in the Broader Impact subsection.
    \item[] Guidelines:
    \begin{itemize}
        \item The answer \answerNA{} means that there is no societal impact of the work performed.
        \item If the authors answer \answerNA{} or \answerNo{}, they should explain why their work has no societal impact or why the paper does not address societal impact.
        \item Examples of negative societal impacts include potential malicious or unintended uses (e.g., disinformation, generating fake profiles, surveillance), fairness considerations, privacy considerations, and security considerations.
        \item If there are negative societal impacts, the authors could also discuss possible mitigation strategies.
    \end{itemize}

\item {\bf Safeguards}
    \item[] Question: Does the paper describe safeguards that have been put in place for responsible release of data or models that have a high risk for misuse (e.g., pre-trained language models, image generators, or scraped datasets)?
    \item[] Answer: \answerNA{}
    \item[] Justification: The model is a positional encoding improvement for Transformers and does not introduce new high-risk capabilities beyond the base architecture.
    \item[] Guidelines:
    \begin{itemize}
        \item The answer \answerNA{} means that the paper poses no such risks.
        \item Released models that have a high risk for misuse or dual-use should be released with necessary safeguards to allow for controlled use of the model.
    \end{itemize}

\item {\bf Licenses for existing assets}
    \item[] Question: Are the creators or original owners of assets (e.g., code, data, models), used in the paper, properly credited and are the license and terms of use explicitly mentioned and properly respected?
    \item[] Answer: \answerYes{}
    \item[] Justification: All baselines and datasets are cited. MovieLens is available under its standard research license; Amazon Beauty data is used per the terms of \citet{ni2019justifying}. No code assets are redistributed.
    \item[] Guidelines:
    \begin{itemize}
        \item The answer \answerNA{} means that the paper does not use existing assets.
        \item The authors should cite the original paper that produced the code package or dataset.
        \item The authors should state which version of the asset is used and, if possible, include a URL.
        \item The name of the license (e.g., CC-BY 4.0) should be included for each asset.
    \end{itemize}

\item {\bf New assets}
    \item[] Question: Are new assets introduced in the paper well documented and is the documentation provided alongside the assets?
    \item[] Answer: \answerNA{}
    \item[] Justification: No new datasets are released in this paper. Code will be released with documentation upon acceptance.
    \item[] Guidelines:
    \begin{itemize}
        \item The answer \answerNA{} means that the paper does not release new assets.
        \item Researchers should communicate the details of the dataset/code/model as part of their submissions via structured templates. This includes details about training, license, limitations, etc.
    \end{itemize}

\item {\bf Crowdsourcing and research with human subjects}
    \item[] Question: For crowdsourcing experiments and research with human subjects, does the paper include the full text of instructions given to participants and screenshots, if applicable, as well as details about compensation (if any)?
    \item[] Answer: \answerNA{}
    \item[] Justification: The paper does not involve crowdsourcing or human subjects experiments.
    \item[] Guidelines:
    \begin{itemize}
        \item The answer \answerNA{} means that the paper does not involve crowdsourcing nor research with human subjects.
        \item According to the NeurIPS Code of Ethics, workers involved in data collection, curation, or other labor should be paid at least the minimum wage in the country of the data collector.
    \end{itemize}

\item {\bf Institutional review board (IRB) approvals or equivalent for research with human subjects}
    \item[] Question: Does the paper describe potential risks incurred by study participants, whether such risks were disclosed to the subjects, and whether Institutional Review Board (IRB) approvals (or an equivalent approval/review based on the requirements of your country or institution) were obtained?
    \item[] Answer: \answerNA{}
    \item[] Justification: No human subjects were involved in this research.
    \item[] Guidelines:
    \begin{itemize}
        \item The answer \answerNA{} means that the paper does not involve crowdsourcing nor research with human subjects.
        \item Depending on the country in which research is conducted, IRB approval (or equivalent) may be required for any human subjects research. If you obtained IRB approval, you should clearly state this in the paper.
    \end{itemize}

\item {\bf Declaration of LLM usage}
    \item[] Question: Does the paper describe the usage of LLMs if it is an important, original, or non-standard component of the core methods in this research? Note that if the LLM is used only for writing, editing, or formatting purposes and does \emph{not} impact the core methodology, scientific rigor, or originality of the research, declaration is not required.
    \item[] Answer: \answerNA{}
    \item[] Justification: LLMs are discussed as a future extension target in Section~\ref{sec:conclusion}, not as a component of the methodology. No LLM was used in the core method development.
    \item[] Guidelines:
    \begin{itemize}
        \item The answer \answerNA{} means that the core method development in this research does not involve LLMs as any important, original, or non-standard components.
        \item Please refer to our LLM policy in the NeurIPS handbook for what should or should not be described.
    \end{itemize}

\end{enumerate}



\end{document}
